\section{Seperation of Concerns}

\begin{frame}{Separation of Concerns}
\begin{block}{Separation of Concerns}
Separation of concerns is a design principle for separating a computer program into distinct sections such that each section addresses a separate concern. A concern is a set of information that affects the code of a computer program.
\end{block}
\end{frame}

\begin{frame}{Separation of Concerns}
\begin{itemize}
    \item One of the most important concepts of OOP! (together with encapsulation).
    \item Classes separate our program into different parts
    \item This makes it so that one part of the program does not have to be \textbf{concerned} with another part of our program.
    \item \textbf{Encapsulation} facilitates this by hiding the implementation of our classes behind well defined interfaces (e.g. getters and setters).
    \begin{itemize}
        \pause
        \item So other classes don't have to be concerned with how a variable is changed or read.
    \end{itemize}
    \item \textbf{In a nutshell:} a class should only do what it is designed to do.
\end{itemize}
\end{frame}

\begin{frame}{Why you should use OOP!}
Why you should know about seperation of concerns and encapsulation.
\end{frame}

\begin{frame}{Why you should use OOP!}
Let's say you have a group project, to create a phone book app.
\begin{itemize}
    \pause
    \item You create a \textbf{PhoneBook} class that stores phone numbers.
    \pause
    \item You want to be able to add and see phone numbers via the terminal.
    \pause
    \item (Following seperation of concerns) you create a different class \textbf{PhoneBookApplication} that handles the input, output, and the loop of your program.
    \pause
    \item You also want to maintain a database with all phone numbers (e.g. txt file, CSV file, or MySQL database).
    \pause
    \item So, you create another class that interacts with a database \textbf{DatabaseHandler}
\end{itemize}
\end{frame}

\begin{frame}{Why you should use OOP!}
\begin{itemize}
    \item Let's assume person A made the PhoneBook class, person B made the PhoneBookApplication class, and person C made the DatabaseHandler class.
    \pause
    \item By creating different classes you can divide the work nicely!
    \pause
    \item Lets assume all classes have well defined methods.
    \pause
    \item For example, the DatabaseHandler has a method save\_to\_database(phonebook).
    \pause
    \item Only person C has to be concerned with how this method works, and has to make sure that it saves to a database.
    \pause
    \item Person B needs to use this method, but does not need to be \textbf{concerned} with how it works! The implementation is \textbf{encapsulated} behind the method.
\end{itemize}
\end{frame}

\begin{frame}{Why you should use OOP!}
    \begin{itemize}
        \item Separation of Concerns and encapsulation makes creating big projects (in groups) way easier!
    \end{itemize}
\end{frame}